\subsection{\texorpdfstring{映射 \((\mathbf{G} / \mathbf{T})\times \mathbf{T}\to \mathbf{G}\) 与 \((\mathbf{G} / \mathbf{T})\times \mathbf{A}\rightarrow \mathbf{G}_r\)}{映射 (\textbackslash mathbf\{G\} / \textbackslash mathbf\{T\})\textbackslash times \textbackslash mathbf\{T\}\textbackslash to \textbackslash mathbf\{G\} 与 (\textbackslash mathbf\{G\} / \textbackslash mathbf\{T\})\textbackslash times \textbackslash mathbf\{A\}\textbackslash rightarrow \textbackslash mathbf\{G\}\_r}}

从 G × T 到 G 的映射 \((g,t)\mapsto gtg^{-1}\) 通过过渡到商诱导出解析流形间的一个态射

\[
f: (\mathrm{G} / \mathrm{T}) \times \mathrm{T} \to \mathrm{G},
\]

它是满射的（§2，第 2 节，定理 2）。通过限制，\(f\) 诱导出一个满射态射

\[
f _ {r}: (\mathrm{G} / \mathrm{T}) \times \mathrm{T} _ {r} \rightarrow \mathrm{G} _ {r}.
\]

与 \(Id_{G/T} \times exp_{T}\) 复合，我们得到同样是满射的态射

\[
\begin{array}{c} \varphi : (\mathrm{G/T}) \times \mathfrak {t} \to \mathrm{G}, \\ \varphi_ {r}: (\mathrm{G/T}) \times \mathfrak {t} _ {r} \to \mathrm{G} _ {r}; \end{array}
\]

最后，若 A 是 t 的一个单房，则 \(\varphi_{r}\) 诱导出一个满射态射

\[
\varphi_ {\mathrm{A}}: (\mathrm{G/T}) \times \mathrm{A} \to \mathrm{G} _ {r}.
\]

我们如下定义 W 在 G/T 上的一个右作用：设 \(w \in W\) 与 \(u \in G/T\)；把 w 提升为 \(\mathrm{N}_{\mathrm{G}}(\mathrm{T})\) 的一个元素 n，把 u 提升为 G 的一个元素 g。则 gn 在 G/T 中的像不依赖于 n 与 g 的选取；我们把它记作 u.w。

对这个作用，W 自由地作用于 G/T 上：实际上，沿用前面的记号，假设 u.w = u；则 \(gn \in gT\)，故 \(n \in T\) 且 w = 1。

我们定义 W 在 \((\mathrm{G}/\mathrm{T}) \times \mathrm{T}\) 上的一个右作用如下：

\[
(u, t). w = (u. w, w ^ {- 1} (t)), \quad u \in \mathrm{G} / \mathrm{T}, t \in \mathrm{T}, w \in \mathrm{W}
\]

再定义 \(W_{a}^{\prime}\) 在 (G/T) × t 上的一个右作用如下：

\[
(u, x). \omega = (u. \overline {{\omega}}, \omega^ {- 1} (x)), \quad u \in \mathrm{G/T}, x \in \mathfrak {t}, \omega \in \mathrm{W} _ {a} ^ {\prime},
\]

其中 \(\overline{\omega}\) 是 \(\omega\) 在商 \(W_{a}^{\prime}/\Gamma(T)=W\) 中的像。

若 A 是 t 的一个单房，且 \(H_{A}\) 是 \(W_{a}^{\prime}\) 中稳定 A 的子群，则通过限制得到 \(H_{A}\) 在 \((\mathrm{G}/\mathrm{T}) \times \mathrm{A}\) 上的一个作用。

这些不同的作用与态射 \(f, \varphi\) 及 \(\varphi_{\mathrm{A}}\) 相容：对 \(u \in \mathrm{G} / \mathrm{T}, t \in \mathrm{T}, x \in \mathfrak{t}, y \in \mathrm{A}, w \in \mathrm{W}, \omega \in \mathrm{W}_a', h \in \mathrm{H}_{\mathrm{A}}\)，我们有

\[
f ((u, t). w) = f (u, t), \varphi ((u, x). \omega) = \varphi (u, x), \varphi_ {\mathrm{A}} ((u, y). h) = \varphi_ {\mathrm{A}} (u, y).
\]

引理 4. 设 \(g \in \mathrm{G}\)，\(t \in \mathrm{T}\)，并设 \(\bar{g}\) 是 \(g\) 在 \(\mathrm{G} / \mathrm{T}\) 中的像。把 \(\mathrm{G} / \mathrm{T}\)（相应地 T，相应地 G）在 \(\bar{g}\)（相应地 t，相应地 \(gtg^{-1}\)）处的切空间与 \(\mathfrak{g} / \mathfrak{t}\)（相应地 \(\mathfrak{t}\)，相应地 \(\mathfrak{g}\)）等同起来，其借助的是左平移 \(\gamma(g)\)（由 \(g\)（相应地 \(t\)，相应地 \(gtg^{-1}\)）所作）。于是，\(f\) 在 \((\bar{g}, t)\) 处的切线性映射就等同于按下述方式定义的线性映射 \(f': (\mathfrak{g} / \mathfrak{t}) \times \mathfrak{t} \to \mathfrak{g}\)：若 \(z \in \mathfrak{g}, x \in \mathfrak{t}\)，且用 \(\bar{z}\) 表示 \(z\) 在 \(\mathfrak{g} / \mathfrak{t}\) 中的像，则

\[
f ^ {\prime} (\bar {z}, x) = (\mathrm{Ad} g t ^ {- 1}) (z - (\mathrm{Ad} t) z + x).
\]

设 \(\mathrm{F}\) 是从 \(\mathrm{G} \times \mathrm{T}\) 到 \(\mathrm{T}\) 的映射，满足 \(\mathrm{F}(g,t) = gtg^{-1}\)。由于 \(\mathrm{F} \circ (\gamma(g), \mathrm{Id}_{\mathrm{T}}) = \operatorname{Int} g \circ \mathrm{F}\)，我们有 \(\mathrm{T}_{(g,t)}(\mathrm{F})(gz, tx) = \mathrm{T}_t(\operatorname{Int} g) \circ \mathrm{T}_{(e,t)}(\mathrm{F})(z, tx)\)；由第三章 §3 第 12 节命题 46，

\[
\mathrm{T} _ {(e, t)} (\mathrm{F}) (z, t x) = t ((\mathrm{Ad} t ^ {- 1}) z - z) + t x = t ((\mathrm{Ad} t ^ {- 1}) (z - (\mathrm{Ad} t) z + x))
\]

从而

\[
\mathrm{T} _ {(g, t)} (\mathrm{F}) (g z, t x) = g t g ^ {- 1} ((\mathrm{Ad}   g t ^ {- 1}) (z - (\mathrm{Ad}   t) z + x)).
\]

对这个公式过渡到商，立即得到引理。

命题 4. a) 设 \(g \in \mathbf{G}, t \in \mathrm{T}, x \in \mathfrak{t}\)，并设 \(\bar{g}\) 是 \(g\) 在 \(\mathrm{G} / \mathrm{T}\) 中的像。下列条件等价：

\begin{enumerate}
\def\labelenumi{(\roman{enumi})}
\tightlist
\item
  \(t \in \mathrm{T}_r\)（相应地 \(x \in \mathfrak{t}_r\)）。
\end{enumerate}

(i bis) 元素 \(f(\bar{g}, t)\)（相应地 \(\varphi(\bar{g}, x)\)）在 G 中是正则的。

\begin{enumerate}
\def\labelenumi{(\roman{enumi})}
\setcounter{enumi}{1}
\tightlist
\item
  映射 \(f\)（相应地 \(\varphi\)）在点 \((\bar{g}, t)\)（相应地 \((\bar{g}, x)\)）处是浸没。
\end{enumerate}

(ii bis) 映射 \(f\)（相应地 \(\varphi\)）在点 \((\bar{g}, t)\)（相应地 \((\bar{g}, x)\)）处是平展的。

\begin{enumerate}
\def\labelenumi{\alph{enumi})}
\setcounter{enumi}{1}
\item
  映射 \(f_{r}\)（相应地 \(\varphi_{r}\)，相应地 \(\varphi_{A}\)）使 \((\mathrm{G}/\mathrm{T}) \times \mathrm{T}_{r}\)（相应地 \((\mathrm{G}/\mathrm{T}) \times \mathfrak{t}_{r}\)，相应地 \((\mathrm{G}/\mathrm{T}) \times \mathrm{A}\)）成为 \(G_{r}\) 的一个以 W（相应地 \(W_{a}^{\prime}\)，相应地 \(H_{A}\)）为群的主覆盖。
\item
  (i) 与 (i bis) 的等价性是明显的；(ii) 与 (ii bis) 的等价性由关系式 \(\dim((\mathrm{G}/\mathrm{T})\times\mathrm{T})=\dim((\mathrm{G}/\mathrm{T})\times\mathfrak{t})=\dim(\mathrm{G})\) 得出。由引理 4，f 在点 \((\bar{g},t)\) 处是浸没当且仅当 \(g=t+Im(Ad t-Id)\)，而这意味着 t 是正则的。最后，由于 \(\varphi=f\circ(Id_{G/T}\times\exp_{T})\)，\(\varphi\) 在点 \((\bar{g},x)\) 处是平展的当且仅当 f 在点 \((\bar{g},\exp x)\) 处是平展的，由前所述，这就是说 x 属于 \(t_{r}\)。
\item
  于是，态射 \(f_{r}, \varphi_{r}, \varphi_{A}\) 都是平展的。另一方面，W 自由地作用于 G/T 上，从而更是自由地作用于 \((\mathrm{G}/\mathrm{T}) \times \mathrm{T}\) 上。设 \(g, g'\)（G 中）与 \(t, t'\)（\(T_{r}\) 中）满足 \(f(\bar{g}, t) = f(\bar{g}', t')\)；则 \(\operatorname{Int} g^{-1} g'\) 把 \(t'\) 映到 t，从而正规化 T，因为 \(\mathrm{T} = \mathrm{Z}(t)_{0} = \mathrm{Z}(t')_{0}\)，且 \(g^{-1} g'\) 在 W 中的类 w 把 \((\bar{g}, t)\) 映到 \((\bar{g}', t')\)。由此可知，\(f_{r}\) 是以 W 为群的主覆盖；这立即蕴含 \(\varphi_{r}\) 是以 \(W_{a}'\) 为群的主覆盖，从而通过限制到 \((\mathrm{G}/\mathrm{T}) \times \mathrm{A}\)（\((\mathrm{G}/\mathrm{T}) \times \mathbf{t}_{r}\) 的连通分支）上，得到 \(\varphi_{A}\) 是以 \(H_{A}\) 为群的主覆盖。
\end{enumerate}

注. 1) 由 §2 第 4 节命题 3，流形 \((\mathrm{G}/\mathrm{T}) \times \mathrm{A}\) 是单连通的。由此可知，\(\varphi_{A}\) 是 \(G_{r}\) 的一个泛覆盖；由于 \(\pi_{1}(G_{r})\) 典则地同构于 \(\pi_{1}(G)\)（第 1 节，命题 1 与《一般拓扑》，

第十一章，编写中），我们重新得到 \(\pi_{1}(\mathrm{G})\) 可等同于 \(H_{A}\)（即等同于 \(\Gamma(\mathrm{T})/\mathrm{N}(\mathrm{G},\mathrm{T})\)）这一事实。

\begin{enumerate}
\def\labelenumi{\arabic{enumi})}
\setcounter{enumi}{1}
\item
  \(\varphi_{A}\) 在 \(W \times A \subset (G/T) \times A\) 上的限制使 \(W \times A\) 成为 \(T_{r}\) 的一个以 \(H_{A}\) 为群的主覆盖。于是我们重新得到第 2 节命题 2 的推论 1。
\item
  用 \(\mathfrak{g}_r\) 表示在指数映射下 \(\mathrm{G}_r\) 的逆像，并用 \(\varepsilon : \mathfrak{g}_r \to \mathrm{G}_r\) 表示由 \(\exp_{\mathrm{G}}\) 诱导的映射。映射 \((g, x) \mapsto (\operatorname{Ad} g)(x)\)（从 \(\mathrm{G} \times \mathfrak{t}_r\) 到 \(\mathfrak{g}_r\)）通过过渡到商定义出一个映射 \(\psi_r : (\mathrm{G} / \mathrm{T}) \times \mathfrak{t}_r \to \mathfrak{g}_r\)。我们有 \(\varepsilon \circ \psi_r = \varphi_r\)。设 \(w \in \mathrm{W}, \gamma \in \Gamma(\mathrm{T})\) 与 \(\omega \in \mathrm{W}_a'\) 满足 \(\omega(z) = w(z) + \gamma\)（对所有 \(z \in \mathfrak{t}\)）；则 \(\psi_r((\bar{g}, x)\omega) = \psi_r(\bar{g}, x) - (\operatorname{Ad} g)(\gamma)\) 对 \(g \in \mathrm{G}, x \in \mathfrak{t}_r\) 成立，故 \(\psi_r((\bar{g}, x)\omega) = \psi_r(\bar{g}, x)\) 当且仅当 \(\gamma = 0\)。由此可知（参看《一般拓扑》，第十一章，编写中），\(\psi_r\) 是 \(\mathfrak{g}_r\) 的一个以 W 为群的主覆盖，而 \(\varepsilon : \mathfrak{g}_r \to \mathrm{G}_r\) 是相伴于主覆盖 \(\varphi_r\) 的一个覆盖，其纤维同构于 \(\mathrm{W}_a'\)-集合 \(\mathrm{W}_a'/\mathrm{W}\)。
\end{enumerate}

